
%\documentclass[10pt,twoside,twocolumn]{article}
\documentclass[12pt,twoside]{article}
\usepackage[bf,small]{caption}
\usepackage[letterpaper,hmargin=1in,vmargin=1in]{geometry}
\usepackage{paralist} % comapctitem, compactdesc, compactenum
\usepackage{titlesec}
\usepackage{titletoc}
\usepackage{times}
\usepackage{hyperref}
\usepackage{algorithmic}
\usepackage{graphicx}
\graphicspath{{./graphics/}}
\usepackage{xspace}
\usepackage{verbatim}
\usepackage{url}
\usepackage{float}
\hyphenation{Sub-Bytes Shift-Rows Mix-Col-umns Add-Round-Key}

\setlength{\parskip}{12pt}
\setlength{\parindent}{0pt}

\newcommand{\hdb}{\emph{hashdb}\xspace}
\newcommand{\bulk}{\emph{bulk\_extractor}\xspace}
\newcommand{\mdd}{\emph{md5deep}\xspace}
\newcommand{\bev}{\emph{Bulk Extractor Viewer}\xspace}

\begin{document}

\begin{center}
\Large Demo: Finding Similarities between Media Images \\
\large Using \bulk and \hdb
\end{center}

In this demo, we find similarities between media images
by finding block hashes that are common between them.
Here are the steps:
\begin{compactenum}
\item Generate a hash database of block hashes from media image 1.
\item Generate a hash database of block hashes from media image 2.
\item Obtain common block hashes
by taking the intersection of these two databases.
\end{compactenum}
Here is the workflow:
\begin{figure}[H]
  \center
  \includegraphics[scale=0.6]{drawings/import_hashes}
  \caption*{Run \bulk with the \hdb \texttt{import} option selected \\
            to import media images into hash databases.}
\end{figure}

\begin{figure}[H]
  \center
  \includegraphics[scale=0.6]{drawings/intersect_hashes}
  \caption*{Run the \hdb \texttt{intersect} command
            to generate the database of common hashes.}
\end{figure}

Media Images contain many false positives because they contain
null blocks, operating system files, common header blocks, etc.
To remove these false positives from our hash dtabase,
we subtract out hashes if they are are already known:

\begin{figure}[H]
  \center
  \includegraphics[scale=0.6]{drawings/subtract_whitelist}
  \caption*{Run the \hdb \texttt{subtract} command to create a database without specified hashes.}
\end{figure}

Now view the list of forensic paths common to both media images:

\begin{figure}[H]
  \center
  \includegraphics[scale=0.6]{drawings/print_sources}
  \caption*{Run the \hdb \texttt{get\_sources} command \\
            to list common hashes and where they were sourced from.}
\end{figure}

%steps
Here are the steps to perform this demo:
\begin{compactenum}
\item Download and install \hdb  and \bulk compiled with the \hdb scanner
from
\url{http://digitalcorpora.org/downloads/hashdb}
as described in the Scanner demo
for finding part of a video file in a media image,
\url{http://digitalcorpora.org/downloads/hashdb/demo/scanner\_demo.pdf}
\item Identify two media images to compare.
I'll call them \texttt{media\_image\_1} and \texttt{media\_image\_2}.
Note that images are available at
\url{http://digitalcorpora.org/corpora/disk-images}.
For example
\url{http://digitalcorpora.org/corp/nps/scenarios/2009-m57-patents/drives-redacted/jo-favorites-usb-2009-12-11.E01}
and
\url{http://digitalcorpora.org/corp/nps/scenarios/2009-m57-patents/drives-redacted/jo-work-usb-2009-12-11.E01}.

\item Import media image 1:
\begin{verbatim}
$ bulk_extractor -o outdir1 -S hashdb_mode=import media_image_1
\end{verbatim}

\item Import media image 2:
\begin{verbatim}
$ bulk_extractor -o outdir2 -S hashdb_mode=import media_image_2
\end{verbatim}

\item Create the intersection of media image 1 and 2
to obtain a database of hashes that are common between them:
\begin{verbatim}
$ hashdb create intersection.hdb
$ hashdb intersect media_image1/hashdb.hdb media_image2/hashdb.hdb \
         intersection.hdb
\end{verbatim}
\end{compactenum}

Now database \texttt{intersection.hdb} contains common hashes,
but we want to know what is in it.
Here are some approaches for further analysis:
\begin{compactitem}
\item There are ``false positives'' from system files.
Which hashes are actually interesting user data?
Subtract false positives from a database of files known to be not interesting,
such as system files,
so that the remaining hashes are pertinent data common between the
media images:
\begin{verbatim}
$ hashdb create intersection2.hdb
$ hashdb subtract intersection.hdb false_positives.hdb \\
         intersection2.hdb
\end{verbatim}

\item What hashes match previously encountered data?
Intersect the media image intersection database
with a database of known interesting data,
and print out the list of files from which the known hashes were sourced:
\begin{verbatim}
$ hashdb create recognized_hashes.hdb
$ hashdb intersect intersection.hdb known_hashes.hdb \
         recognized_hashes.hdb
$ hashdb get_sources recognized_hashes.hdb > output.txt
\end{verbatim}
\end{compactitem}

Note that these approaches for further analysis require
databaes of previously encountered files.
For a demo on preparing a block hash database, please see
\url{http://digitalcorpora.org/downloads/hashdb/demo/create\_hdb\_demo.pdf}.
To create a database of ``false positives'',
import hashes from some media images or from the NSRL.

Note that this demo uses the default hash block size of 4096 bytes
and the default sector size of 512 bytes.
This is not required.
For example to perform a similarity analysis
with a block fragment resolution of 64 bytes along every byte boundary,
use a block size of 64 bytes and a sector size of 1 byte.
The syntax for importing media images 1 and 2 would be:
\begin{verbatim}
$ bulk_extractor -o outdir1 \
                 -S hashdb_mode=import \
                 -S hashdb_block_size=64 \
                 -S hashdb_sector_size=1 \
                 media_image_1
$ bulk_extractor -o outdir2 \
                 -S hashdb_mode=import \
                 -S hashdb_block_size=64 \
                 -S hashdb_sector_size=1 \
                 media_image_2
\end{verbatim}

\end{document}

